\begin{table}[H]
\centering
\caption{Optimal hedge ratio for amerikanske aktiver før og efter april 2025}
\label{tab:campbell_brud}
\begin{threeparttable}
\begin{tabular}{l cccc c}
\toprule
 & (1) & (2) & (3) & (4) & \\
 & Hele perioden & Før april 2025 & Efter april 2025 & Forskel & $n_{\text{post}}$ \\
\midrule
\multicolumn{6}{l}{\textit{Panel A: Kvartalsafkast (baseline)}} \\
Aktier & $0.49^{**}$ & $0.35^{***}$ & $1.47^{*}$ & $1.12^{***}$ & 18 \\
 & $(0.25)$ & $(0.23)$ & $(0.25)$ & $(0.32)$ & \\[2pt]
Obligationer & $0.75^{***}$ & $0.72^{***}$ & $0.86^{*}$ & $0.14$ & 18 \\
 & $(0.08)$ & $(0.09)$ & $(0.08)$ & $(0.12)$ & \\[2pt]
\addlinespace
\multicolumn{6}{l}{\textit{Panel B: Månedsafkast}} \\
Aktier & $0.46^{***}$ & $0.45^{***}$ & $0.55$ & $0.10$ & 18 \\
 & $(0.18)$ & $(0.20)$ & $(0.49)$ & $(0.53)$ & \\[2pt]
Obligationer & $0.77^{***}$ & $0.77^{***}$ & $0.77^{**}$ & $-0.00$ & 18 \\
 & $(0.05)$ & $(0.06)$ & $(0.11)$ & $(0.12)$ & \\[2pt]
\addlinespace
\multicolumn{6}{l}{\textit{Panel C: Ugeafkast}} \\
Aktier & $0.52^{***}$ & $0.52^{***}$ & $0.47$ & $-0.05$ & 78 \\
 & $(0.13)$ & $(0.14)$ & $(0.40)$ & $(0.42)$ & \\[2pt]
Obligationer & $0.90^{***}$ & $0.89^{***}$ & $1.00$ & $0.12$ & 78 \\
 & $(0.03)$ & $(0.03)$ & $(0.19)$ & $(0.19)$ & \\[2pt]
\bottomrule
\end{tabular}
\begin{tablenotes}[flushleft]\footnotesize
\item \textit{Note:} Optimal hedge ratio $h^* = 1 + \beta$, hvor $\beta$ er koefficienten fra en regression af merafkastet på amerikanske aktier (S\&P 500 Total Return) og statsobligationer (Bloomberg US Treasury) på dollarens merafkast over for euroen (Campbell et al., 2010). (1) Hele stikprøven. (2) $h^*$ før april 2025, $1+\beta$, og (3) efter april 2025, $1+\beta+\phi$, begge fra en model med brud-dummy. (4) Forskellen $\phi$. Standardfejl i parentes: Newey-West med 6 lags i panel A (overlappende kvartalsafkast) og HC3 i panel B og C. Stjerner ved $h^*$ angiver, at $h^*$ er signifikant forskellig fra 1, svarende til Campbell et al.'s (2010) test af, at den optimale dollareksponering, $1-h^*$, er nul. Stjerner ved forskellen angiver, at $h^*$ har ændret sig. $^{*}$, $^{**}$ og $^{***}$ angiver signifikans på 10, 5 og 1 pct.-niveau. $n_{\text{post}}$ er antallet af observationer efter bruddet (måneder i panel A og B, uger i panel C); i panel A overlapper kvartalsafkastene, så antallet af uafhængige kvartaler er ca. en tredjedel.
\end{tablenotes}
\end{threeparttable}
\end{table}